% Target: rmp-iii-b-braid-one
% Formula: The first braid panel moves string a past anyon i and string b.
% Ink: Canonical public kernel plane patch with two crossing strings.
% Boundary: Eight lattice legs run past the patch and four site legs leave the
%   plane; both strings end below the lower rail.
% Source: paper TN-Review-main.tex line 1675; author source ImagesReview Section
%   III/ImagesReview Section III.tex lines 799-825
% Capabilities: kernel, plane-frame, multi-strand-braid, crossing-order
\[
  \tenkzkernel
  \tndeclare{species}{operator}{hue=source:red}
  % The author draws all four lattice lines inside one plane, so every one of
  % them carries a virtual bond; the site index is the stub that leaves the
  % plane, and the four stubs end in air.
  \begin{tenkz}[lattice={2x2}, frame=plane, physical=up,
                west=open, east=open, north=open, south=open]
    % the two-by-two lattice patch: anyons j and i on the near rail
    \tn[at=(1,1), name=j, label pos=se]{j}
    \tn[at=(1,2), name=i, label pos=se]{i}
    \tn[at=(2,1), name=p]{}
    \tn[at=(2,2), name=q]{}
    % string b leaves anyon j, laps over the lattice line that runs north out
    % of j, and exits below; the declared crossing orders string over bond, as
    % the source paints it.  The wire is a junction chain (retired via=): an
    % invisible junction at each of the two waypoints with no marker bead of
    % its own (jb1, jb2), then the chain runs through the marker beads bt, bs
    % that already sit at the remaining two waypoints, so no atom is
    % duplicated there.  Segment b2 is the one that laps over the lattice
    % line, so it carries the renamed crossing declaration; every consecutive
    % pair of segments also needs its own declared order at the corner they
    % share, same as the torus-two wrap chain.
    \tn[at=0.3 ne of j, skin=none, name=jb1]{}
    \tn[at=0.3 nw of j, skin=none, name=jb2]{}
    \tn[at=0.4 w of j, name=bt]{}
    \tn[at=0.4 w of p, name=bs]{}
    \tnwire[kind=string, species=operator, name=b1, route=arc]{0.1 e of j}{jb1}
    \tnwire[kind=string, species=operator, name=b2, route=arc,
            cross={over at crossing of b1 and b2,
                   over at crossing of b2 and open-north-col-1}]{jb1}{jb2}
    \tnwire[kind=string, species=operator, name=b3, route=arc,
            cross={over at crossing of b2 and b3}]{jb2}{bt}
    \tnwire[kind=string, species=operator, name=b4, route=arc,
            cross={over at crossing of b3 and b4}]{bt}{bs}
    \tnwire[kind=string, species=operator, name=b5, route=arc,
            cross={over at crossing of b4 and b5}]{bs}{open s}
    % marker bead where string b pierces the lattice line, and label bead for
    % string a
    \tn[at=crossing of b2 and open-north-col-1, name=bb, label pos=e]{b}
    \tn[at=midway p and q, name=am, label pos=se]{a}
    % string a rises from below through the lower rail into anyon i; its
    % first waypoint is the marker bead am already declared above, its
    % second waypoint has no bead so it gets an invisible junction (ja2)
    \tn[at=0.4 sw of i, skin=none, name=ja2]{}
    \tnwire[kind=string, species=operator, name=a1, route=arc]{open s}{am}
    \tnwire[kind=string, species=operator, name=a2, route=arc,
            cross={over at crossing of a1 and a2}]{am}{ja2}
    \tnwire[kind=string, species=operator, name=a3, route=arc,
            cross={over at crossing of a2 and a3}]{ja2}{i}
  \end{tenkz}
\]
